\clearpage
\appendix
\label{appendix}

\section{Experiment Setup}
\label{sec:app-experiment}
\subsection{Datasets} 
We select four popular complex KGQA datasets as in-domain datasets, \ie \textit{WebQuestionsSP (WebQSP)}~\cite{WebQSP}, \textit{Complex WebQuestions 1.1 (CWQ)}~\citep{CWQ}, and \textit{GrailQA}~\citep{GrailQA}, which are based on Freebase, and \textit{KQA Pro}~\cite{KQA-Pro}, which is based on Wikidata.
And we select three representative ODQA datasets as out-domain datasets, which are \textit{WebQuestions~(WQ)}~\cite{Berant-WQ-ACL-13}, \textit{Natural Questions~(NQ)}~\cite{Chen-NQ-ACL-17}, and \textit{TriviaQA~(TQ)}~\cite{Joshi-TQ-ACL-17}. Since we only rely on the KG to answer questions, we filter the questions in ODQA datasets that can not be linked to any entity in KG, denoted as \textit{WQ-Freebase}, \textit{NQ-Wiki}, and \textit{TQ-Wiki}, respectively.
Besides, we further select the \textit{MetaQA}~\cite{zhang-aaai-18}, which is based on a domain-specific movie KG, to evaluate the generalibility of our method.
The detail description of these selected datasets is as follows:

$\bullet$ \textbf{WebQSP} consists of 4,737 questions. The answer entities are within a maximum of 2 hops from the topic entity on the Freebase KG. We adopt the train/valid/test splits from GraftNet~\citep{GraftNet} for consistency.

$\bullet$ \textbf{CWQ} is constructed based on WebQSP, which is more challenging. It complicates WebQSP by extending the question entities or adding constraints to restrict the answers. The answer entities are within a maximum of 4 hops from the topic entity on the Freebase KG.

$\bullet$ \textbf{GrailQA} consists of 64,331 questions. Compared to WebQSP and CWQ, it focuses on a more comprehensive generalization capability evaluation from three levels~(\ie i.i.d, compositional, and zero-shot).

$\bullet$ \textbf{KQA Pro} consists of 117,970 questions. The above three datasets are based on Freebase, and it is based on Wikidata, and require multiple reasoning capabilities including compositional reasoning, multi-hop reasoning, quantitative comparison, set operations, and etc.

$\bullet$ \textbf{MetaQA} comprises over 400,000 questions based on a movie domain KG, with answer entities located up to three hops away from
the topic entities. Based on the number of hops, the dataset is divided into three sub-datasets: MetaQA-1hop, MetaQA-2hop, and MetaQA-3hop. Following existing work~\cite{NSM}, we randomly sample just one training case for each question template from the original training set, to form a one-shot training dataset.

$\bullet$ \textbf{WQ} consists of 6,642 questions. The questions are mostly centered around a single named entity and are supposed to be answerable by Freebase KG. We extract xx questions from the original test set to compose the WQ-freebase subset.

$\bullet$ \textbf{NQ} consists of 323,045 questions. Each example contains a question from the Google search and the corresponding answers, which are text spans on the Wikipedia page. Following existing work~\cite{pack-knowledge}, we use the open version of this dataset which discards answers with more than 5 tokens. We extract xx questions from the original test set to compose the NQ-Wiki subset.

$\bullet$ \textbf{TQ} consists of 110K questions. Each example contains a question authored by trivia enthusiasts, and the answers are text spans from the Web or Wikipedia. Following existing work~\cite{pack-knowledge}, we use its unfiltered version for evaluation. We extract xx questions from the original test set to compose the TQ-Wiki subset.

\subsection{Evaluation Protocol} For KGQA, following existing work~\cite{GraftNet}, we use Hits@1 and F1 metrics for WebQSP and CWQ datasets, F1 metric for GrailQA dataset, and Hits@1 for MetaQA. The Hits@1 evaluates the correctness of the top-ranked answer while F1 considers coverage of all the predicted answers.
It's worth noting that some baselines and our approach would return all the unordered answers at the end, which is not suitable for the Hist@1 metric. For a comprehensive comparison, we randomly select one answer per question as the top-ranked answer and then calculate the average Hits@1 result by repeating this process 100 times following existing work~\cite{TIARA}.
For ODQA, following existing work~\cite{pack-knowledge}, we report the EM metric, which evaluates whether the predicted answer is the same as the gold one after performing normalization.

\subsection{Baselines for Comparison} 
For KGQA, we consider the following three types of baseline methods for performance comparison:

$\bullet$ \textbf{subgraph-based reasoning} methods which perform answer reasoning in a retrieval subgraph form KG, including GrafeNet~\cite{GraftNet}, NSM~\cite{NSM}, SubgraphRetrieval~\cite{Zhang-ACL-22}, UniKGQA~\cite{jiang-arxiv-22}, and ReasoningLM~\cite{ReasoningLM} for datasets on Freebase, and KVMemNet~\cite{miller-EMNLP-16}, EmbedKGQA~\cite{Saxena-ACL-20}, and RGCN~\cite{RGCN} for datasets on Wikidata;

$\bullet$ \textbf{LM-based seq2seq generation} methods which generate the final SPARQL query by fine-tuning a sequence-to-sequence language model, including RNG-KBQA~\cite{RNG-KBQA}, Uni-Parser~\cite{UniParser}, ArcaneQA~\cite{ArcaneQA}, PanGu w/ T5-3B~\cite{Gu-ACL-23}, TIARA~\cite{TIARA}, and FC-KBQA~\cite{FC-KBQA} for datasets on Freebase, and RNN SPARQL and BART SPARQL~\cite{KQA-Pro} for datasets on Wikidata;

$\bullet$ \textbf{LLM-based} methods which utilize the powerful zero-shot or few-shot capabilities of LLMs to answer the question without fine-tuning, including ROG~\cite{Luo-arxiv-23-Reasoning}, StructGPT~\cite{Jiang-2023-arxiv-StructGPT}, gpt-3.5-turbo-instruct~(Davinvi-003)~\footnote{https://platform.openai.com/docs}, gpt-3.5-turbo~(ChatGPT)~\footnote{https://platform.openai.com/docs}, and gpt-4~(GPT-4)~\footnote{https://platform.openai.com/docs} for both in-domain datasets.

For ODQA, we focus on the closed-book setting where no documents are provided and consider the following two types of baseline methods:

$\bullet$ \textbf{Fine-tune based} methods which learn to predict the answers, including T5-Base, T5-Large, BART-base, and BART-Large from ~\citep{pack-knowledge};

$\bullet$ \textbf{LLM-based} methods which directly answer the questions in zero-shot setting, including gpt-3.5-turbo-instruct~(Davinvi-003) and gpt-3.5-turbo~(ChatGPT).

\subsection{Implementation Details}
For instruction tuning data construction, we randomly sample a total of 10,000 training data from in-domain datasets in a ratio of 1:5:5:10 for WebQSP, KQA Pro, GrailQA, and CWQ according to some prior empirical studies.
Since we focus on the reasoning process over KG, we suppose the entities have been given for each question following existing work~\cite{GraftNet, NSM, Jiang-2023-arxiv-StructGPT}.
For instruction tuning, we use the LLaMA2-7B~\cite{LLaMA2} as our backbone LLM. We use a cosine learning rate schedule with an initial learning rate of \text{2e-5}, a weight decay of 0.1, a batch size of 256, a maximum length of 1500, and finally fine-tune the model for 3 epochs. 
For the relation retrieval model and entity disambiguation model in the semantic tool, we build them following the existing work~\cite{Zhang-ACL-22, TIARA}.

After instruction tuning, for in-domain datasets, we evaluate the performance of our KG-Agent on the test set of CWQ, WebQSP, KQA Pro, and the dev set of GrailQA.
For out-domain datasets, we evaluate the zero-shot performance of our KG-Agent on the NQ-Wiki, TQ-Wiki, and WQ-Freebase.
For the domain specific dataset, \ie MetaQA, we follow existing work~\cite{NSM, Jiang-2023-arxiv-StructGPT} to extract the one-shot tuning subset from the original training set and fine-tune our KG-Agent with it.
When evaluating the performance of Davinci-003, ChatGPT, and GPT4, we use the latest February version of APIs from OpenAI. And for in-domain datasets, we provide six demonstrations for each test question and parse the prediction results following existing work~\cite{ToG,Jiang-2023-arxiv-StructGPT}, we show the prompt with demonstration for each dataset in Table~\ref{tab:icl_example}.
For the selection of demonstrations, we randomly sample from the corresponding training set for each dataset.
For out-domain datasets, since they are open-domain question answering tasks, we directly input the question to LLMs with proper prompt, as shown in Table~\ref{tab:icl_example}.




\begin{table*}[t]
    \small
    \centering
    \begin{tabular}{c | l}
    \toprule
    \textbf{Dataset} & \textbf{Prompt} \\
    \midrule
    \multirow{20}{*}{\textbf{WebQSP}}  & Question: where is the syracuse university? \\
    & Answer: [New York | Syracuse | United States of America].\\
    & Question: where is the mtv headquarters?\\
    & Answer: [New York City].\\
    & Question: what are the 3 official languages of spain?\\
    & Answer: [Spanish Language].\\
    & Question: what timezone is new england usa in?\\
    & Answer: [Eastern Time Zone].\\
    & Question: who started southwest airlines?\\
    & Answer: [Herb Kelleher | Rollin King].\\
    & Question: what was irving langmuir famous for?\\
    & Answer: [Scientist].\\
    & Question: \{test question\} \\
    & Answer:\\
    \midrule
    \multirow{20}{*}{\textbf{CWQ}}  & Question: Who is the president in the place where the government of Peru is located? \\
    & Answer: [Ollanta Humala].\\
    & Question: Where did Martin Luther King attend university, that has less than 2,586 undergraduates?\\
    & Answer: [Morehouse College].\\
    & Question: What movie produced by the company New Line Cinema was Taylor Lautner in?\\
    & Answer: [Valentine's Day].\\
    & Question: Which year did the team that plays at Turner Field win the World Series?\\
    & Answer: [1995 World Series].\\
    & Question: Which airports are in the circulation area of Il Manifesto?\\
    & Answer: [Leonardo da Vinci–Fiumicino Airport | Ciampino–G. B. Pastine International Airport].\\
    & Question: What were the professions held by the publisher of "The Awakening?"?\\
    & Answer: [Businessperson | Novelist | Writer | Author].\\
    & Question: \{test question\} \\
    & Answer:\\
    \midrule
    \multirow{20}{*}{\textbf{GrailQA}}  & Question: what does the thiokol rocket do? \\
    & Answer: [Launch vehicle].\\
    & Question: what is the club interest of inverness yacht club?\\
    & Answer: [Sailing].\\
    & Question: who is the tour operator of kiribati?\\
    & Answer: [Fly Water Adventures | Kiribati Holidays | Otintaai Tours | Molloy's Tours].\\
    & Question: 1998 marsala vergine terre arse contains what type of grapes?\\
    & Answer: [Catarratto | Grillo | Ansonica].\\
    & Question: \makecell[l]{how many ice hockey coaches have coached the team \\that is currently coached by the eisbaren berlin?}\\
    & Answer: [1].\\
    & Question: court of appeal of sri lanka has what inferior court?\\
    & Answer: [Supreme Court of Sri Lanka].\\
    & Question: \{test question\} \\
    & Answer:\\
    \midrule
    \multirow{20}{*}{\textbf{KQA Pro}}  & Question: Which website officially represents Morgan Creek Productions? \\
    & Answer: [http://www.morgancreek.com/].\\
    & Question: \makecell[l]{Which is shorter: The Killers, with a story set in Los Angeles, \\or Sherlock Holmes, produced by 20th Century Fox?}\\
    & Answer: [Sherlock Holmes].\\
    & Question: What is the street address for the University of San Diego?\\
    & Answer: [5998 Alcala Park, San Diego, CA, 92110-2492].\\
    & Question: How is the Francis Bacon who died in New Haven related to the Yale School of Medicine?\\
    & Answer: [educated at].\\
    & Question: For the film titled Aladdin, where is it published on its publication date of 2019-05-24?\\
    & Answer: [United States of America].\\
    & Question: Who wrote The Postman which was published in 1985?\\
    & Answer: [David Brin].\\
    & Question: \{test question\} \\
    & Answer:\\
    \midrule
    \textbf{\makecell{NQ-Wiki \\ TQ-Wiki \\WQ-Freebase}} & Answer the following question with one or few words. Question: \{test question\}\\
    \bottomrule
    \end{tabular}
    \caption{The prompts used for each dataset when evaluating the ChatGPT, Davinci-003, and GPT-4 models. When performing evaluation, just replace the ``\{test question\}'' with the test question.}
    \label{tab:icl_example}
\end{table*}

\begin{table*}[t]
    \small
    \centering
    \begin{tabular}{c l c}
    \toprule
    \textbf{Type} & \textbf{Tool}  & \textbf{Description}\\
    \midrule
    \multirow{18}{*}{\makecell{Extraction \\ Tool}} & \multirow{2}{*}{get$\_$relation} & Input: entity set $\{e\}$ $\rightarrow$ Output: one-hop relations $R_{\{e\}}$ \\
    & & Return the incoming and outgoing relations of the given entity set $\{e\}$ on KG. \\
    \cmidrule[0.001pt]{2-3}
    & \multirow{2}{*}{get$\_$head$\_$entity} & Input: entity set $\{e\}$, relation $r$ $\rightarrow$ Output: entity set $\{e\}$ \\
    & & Return the head entity set of the given tail entity set $\{e\}$ along the relation $r$. \\
    \cmidrule{2-3}
    & \multirow{2}{*}{get$\_$tail$\_$entity} & Input: entity set $\{e\}$, relation $r$ $\rightarrow$ Output: entity set $\{e\}$ \\
    & & Return the tail entity set of the given head entity set $\{e\}$ along the relation $r$. \\
    \cmidrule{2-3}
    & \multirow{2}{*}{get$\_$entity$\_$by$\_$type} & Input: string type $t$ $\rightarrow$ Output: entity set $\{e\}$ \\
    & & Return the entity set belonging to the given type $t$. \\
    \cmidrule{2-3}
    & \multirow{6}{*}{get$\_$entity$\_$by$\_$constraint} & Input: entity set $\{e\}$, relation $r$, operator $o$, string value $v$ $\rightarrow$ Output: entity set $\{e\}$ \\
    & & \makecell[c]{Return the new entity set whose tail entity along $r$ satisfies the constraint condition. \\If $v$ is not empty, the $o$ should be one of \{``='',``>'',``>='',``<'',``<=''\}, which means \\the comparison between the tail entity and string value should satisfy the operator. \\Else, the $o$ should be one of \{``argmax'',``argmin''\}, which means the tail entity \\should be the maximum or minimum value.}\\
    \cmidrule{2-3}
    & \multirow{2}{*}{get$\_$candidate$\_$entity} & Input: string entity mention $m$ $\rightarrow$ Output: entity set $\{e\}$ \\
    & & Return the candidate linked entity set on the KG for the given entity mention $m$. \\
    \midrule
    \multirow{13}{*}{\makecell{Logic \\ Tool}} & \multirow{2}{*}{count} & Input: entity set $\{e\}$ $\rightarrow$ Output: integer \\
    & & Return the number of entities in the given entity set $\{e\}$. \\
    \cmidrule{2-3}
    & \multirow{2}{*}{intersect} & Input: entity set list $[\{e\}]$ $\rightarrow$ Output: entity set $\{e\}$ \\
    & & Return the intersection of the given list of entity sets.  \\
    \cmidrule{2-3}
    & \multirow{2}{*}{union} & Input: entity set list $[\{e\}]$ $\rightarrow$ Output: entity set $\{e\}$ \\
    & & Return the union of the given list of entity sets.  \\
    \cmidrule{2-3}
    & \multirow{2}{*}{judge} & Input: entity set $\{e\}$, relation $r$, operator $o$, string value $v$ $\rightarrow$ Output: boolean \\
    & & \makecell[c]{Return a boolean value indicating whether the comparison between the tail entity of \\the given entity set $\{e\}$ along relation $r$ and the given value $v$ satisfies the operator $o$.} \\
    \cmidrule{2-3}
    & \multirow{2}{*}{end} & Input: entity set $\{e\}$ $\rightarrow$ Output: entity set $\{e\}$ \\
    & & Return the entity set as the final answer and end the reasoning process. \\
    \midrule
    \multirow{7}{*}{\makecell{Semantic \\ Tool}} & \multirow{3}{*}{retrieve\_relation} & Input: relation set $\{r\}$ $\rightarrow$ Output: relation set $\{r\}$ \\
    & & \makecell[c]{Retrieve relations from the given relation set $\{r\}$ that are \\semantically relevant to the question through neural network.}  \\
    \cmidrule{2-3}
    & \multirow{2}{*}{disambiguate\_entity} & Input: entity set $\{e\}$ $\rightarrow$ Output: entity $e$ \\
    & & \makecell[c]{Disambiguate the candidate linked entity $\{e\}$ based on the question semantics \\and entity information on KG~(\eg one-hop relations) through neural network.} \\
    \bottomrule
    \end{tabular}
    \caption{The detailed definition and usage of all the tools.}
    \label{tab:toolbox}
\end{table*}

\section{Summary of Toolbox}
\label{sec:app-tool}
We summarize the tool name, tool description, and the input argument and output of tools in Table~\ref{tab:toolbox}.